\subsection{\texorpdfstring{映射的计算 \(\operatorname{Tor}^{\mathbf{A}}(g, f)\) 与 \(\operatorname{Ext}_{\mathbf{A}}(f, h)\)}{映射的计算 \textbackslash operatorname\{Tor\}\^{}\{\textbackslash mathbf\{A\}\}(g, f) 与 \textbackslash operatorname\{Ext\}\_\{\textbackslash mathbf\{A\}\}(f, h)}}

设 \(f: \mathbf{M} \to \mathbf{M}'\) , \(h: \mathbf{N}' \to \mathbf{N}\) 分别为左 A-模的同态， \(g: \mathbf{P} \to \mathbf{P}'\) 为右 A-模的同态， \(a: \mathbf{R} \to \mathbf{M}\) , \(a': \mathbf{R}' \to \mathbf{M}'\) , \(b: \mathbf{S} \to \mathbf{P}\) , \(b': \mathbf{S}' \to \mathbf{P}'\) 分别为 M、M'、P、P' 的左分解， \(c: \mathbf{N} \to \mathbf{E}\) , \(c': \mathbf{N}' \to \mathbf{E}'\) 分别为 N 与 N' 的右分解， \(\tilde{f}: \mathbf{R} \to \mathbf{R}'\) , \(\tilde{g}: \mathbf{S} \to \mathbf{S}'\) , \(\tilde{h}: \mathbf{E}' \to \mathbf{E}\) 为复形态射，使得

\[
a ^ {\prime} \circ \tilde {f} = f \circ a, b ^ {\prime} \circ \tilde {g} = g \circ b, \tilde {h} \circ c ^ {\prime} = c \circ h.
\]

命题2. --- 以下两个图表是可交换的：

\[
\begin{array}{c} \operatorname{Tor} ^ {\mathbf {A}} (\mathrm{P}, \mathrm{M}) \xrightarrow {\psi (\mathrm{S} , \mathrm{R})} \mathrm{H} (\mathrm{S} \otimes_ {\mathbf {A}} \mathrm{R}) \\ \operatorname{Tor} ^ {\mathbf {A}} (g, f) \Big \downarrow \qquad \qquad \qquad \qquad \qquad \qquad \qquad \qquad \qquad \qquad \qquad \qquad \qquad \qquad \qquad \qquad \qquad \qquad \qquad \qquad \qquad \qquad \qquad \qquad \qquad \qquad \qquad \qquad \qquad \qquad \qquad \qquad \qquad \qquad \\ \operatorname{Tor} ^ {\mathbf {A}} (\mathrm{P} ^ {\prime}, \mathrm{M} ^ {\prime}) \xrightarrow {\psi (\mathrm{S} ^ {\prime} , \mathrm{R} ^ {\prime})} \mathrm{H} (\mathrm{S} ^ {\prime} \otimes_ {\mathbf {A}} \mathrm{R} ^ {\prime})  , \end{array}
\]

\[
\begin{array}{c} \mathrm{H} (\operatorname{Homgr} _ {\mathbf {A}} (\mathrm{R} ^ {\prime}, \mathrm{E} ^ {\prime})) \xrightarrow {\varphi (\mathrm{R} ^ {\prime} , \mathrm{E} ^ {\prime})} \operatorname{Ext} _ {\mathbf {A}} (\mathrm{M} ^ {\prime}, \mathrm{N} ^ {\prime}) \\ \mathrm{H} (\operatorname{Homgr} _ {\mathbf {A}} (\tilde {f}, \tilde {h})) \Biggl \downarrow \\ \mathrm{H} (\operatorname{Homgr} _ {\mathbf {A}} (\mathrm{R}, \mathrm{E})) \xrightarrow {\varphi (\mathrm{R} , \mathrm{E})} \operatorname{Ext} _ {\mathbf {A}} (\mathrm{M}, \mathrm{N}). \end{array}
\]

设 \(\alpha: L(M) \to R\) , \(\alpha': L(M') \to R'\) , \(\gamma: L(P) \to S\) , \(\gamma': L(P') \to S'\) 为复形态射，使得

\[
a \circ \alpha = p _ {\mathrm{M}}, a ^ {\prime} \circ \alpha^ {\prime} = p _ {\mathrm {M^ {\prime}}}, b \circ \gamma = p _ {\mathrm{P}}, b ^ {\prime} \circ \gamma^ {\prime} = p _ {\mathrm {P^ {\prime}}}.
\]

根据定义， \(\mathrm{H}(\tilde{g}\otimes\tilde{f})\circ\psi(\mathrm{S},\mathrm{R})\) 等于

\[
\mathrm{H} (\tilde {g} \otimes \tilde {f}) \circ \mathrm{H} (\gamma \otimes \alpha) = \mathrm{H} ((\tilde {g} \circ \gamma) \otimes (\tilde {f} \circ \alpha)),
\]

而 \(\psi(\mathbf{S}',\mathbf{R}')\circ\mathrm{Tor}^{\mathbf{A}}(g,f)\) 等于

\[
\mathrm{H} \left(\gamma^ {\prime} \otimes \alpha^ {\prime}\right) \circ \mathrm{H} \left(\mathrm{L} (g) \otimes \mathrm{L} (f)\right) = \mathrm{H} \left(\left(\gamma^ {\prime} \circ \mathrm{L} (g)\right) \otimes \left(\alpha^ {\prime} \circ \mathrm{L} (f)\right)\right).
\]

另一方面， \(\alpha' \circ L(f)\) 与 \(\tilde{f} \circ \alpha\) 是从 L(M) 到 R′ 的两个态射，使得 \(a' \circ (\alpha' \circ L(f)) = p_{\mathrm{M}} \circ L(f) = f \circ p_{\mathrm{M}} = f \circ a \circ \alpha = a' \circ (\tilde{f} \circ \alpha)\) 。根据 X，第 49 页，命题 3， \(\alpha' \circ L(f)\) 与 \(\tilde{f} \circ \alpha\) 是同伦的；同样地 \(\gamma' \circ L(g)\) 与 \(\tilde{g} \circ \gamma\) 是同伦的，因此 \((\gamma' \circ L(g)) \otimes (\alpha' \circ L(f))\) 与 \((\tilde{g} \circ \gamma) \otimes (\tilde{f} \circ \alpha)\) 根据 X，第 64 页，命题 3 也是同伦的。因此我们有

\[
\begin{array}{r l} \mathrm{H} (\tilde {g} \otimes \tilde {f}) \circ \psi (\mathrm{S}, \mathrm{R}) & = \mathrm{H} ((\tilde {g} \circ \gamma) \otimes (\tilde {f} \circ \alpha)) = \mathrm{H} ((\gamma^ {\prime} \circ \mathrm{L} (g)) \otimes (\alpha^ {\prime} \circ \mathrm{L} (f))) \\ & = \psi (\mathrm{S} ^ {\prime}, \mathrm{R} ^ {\prime}) \circ \mathrm{Tor} ^ {\mathrm{A}} (g, f). \end{array}
\]

我们类似地论证第二个图表。

注记. --- 类似地考虑复形态射的交换图表。

\[
\begin{array}{c c c} \mathrm{R} _ {1} \xrightarrow {a _ {1}} \mathrm{M} & \mathrm{S} _ {1} \xrightarrow {b _ {1}} \mathrm{P} & \mathrm{N} ^ {\prime} \xrightarrow {c _ {1} ^ {\prime}} \mathrm{E} ^ {\prime} \\ \tilde {f} \Big \downarrow & f \Big \downarrow & g \Big \downarrow \\ \mathrm{R} _ {1} ^ {\prime} \xrightarrow {a _ {1} ^ {\prime}} \mathrm{M} ^ {\prime} & \mathrm{S} _ {1} ^ {\prime} \xrightarrow {b _ {1} ^ {\prime}} \mathrm{P} ^ {\prime} & \mathrm{N} \xrightarrow {c _ {1}} \mathrm{E} _ {1} \end{array}
\]

其中复形 \(R_{1}\) , \(R_{1}^{\prime}\) , \(S_{1}\) , \(S_{1}^{\prime}\) 是投射的且在右侧为零，而复形 \(E_{1}\) , \(E_{1}^{\prime}\) 是内射的且在左侧为零。那么

\[
\begin{array}{r l} & {\mathrm{Tor} ^ {\mathbf {A}} (g, f) \circ \psi^ {\prime} (\mathrm{S} _ {1}, \mathrm{R} _ {1}) = \psi^ {\prime} (\mathrm{S} _ {1} ^ {\prime}, \mathrm{R} _ {1} ^ {\prime}) \circ \mathrm{H} (\tilde {g} \otimes \tilde {f})} \\ & {\varphi^ {\prime} (\mathrm{R} _ {1}, \mathrm{E} _ {1}) \circ \mathrm{Ext} _ {\mathbf {A}} (f, h) = \mathrm{H} (\mathrm{Homgr} _ {\mathbf {A}} (\tilde {f}, \tilde {h})) \circ \varphi^ {\prime} (\mathrm{R} _ {1} ^ {\prime}, \mathrm{E} _ {1} ^ {\prime})} \end{array}
\]

这由与命题 2 类似的论证得到。

